\begin{frame}{Why no convergence?}
  \begin{itemize}
    \item Difficulties:
    \begin{itemize}
      \item too many possible actions at each step
      \item ideal opponent that never makes mistakes
      \item very sparse positive rewards in training episodes
    \end{itemize}
    \item Possible solutions:
    \begin{itemize}
      \item more training time 
      \item implement draw by repetition and 50-move rule $\rightarrow$ more frequent terminal states
    \end{itemize}
  \end{itemize}
\end{frame}

\begin{frame}{What did we learn?}
  \begin{itemize}
    \item The importance of parameters choice (e.g., rewards)
    \item The danger of introducing a maximum length for each episode
    \item The importance of starting from simpler solutions and then increase complexity
  \end{itemize}
\end{frame}

\begin{frame}
  Why is it so difficult to solve chess endgames?

\begin{quote}
A player can sometimes afford the luxury of an inaccurate move, or even a definite error, in the opening or middlegame without necessarily obtaining a lost position. In the endgame … an error can be decisive, and we are rarely presented with a second chance.
\end{quote}

\hfill - Paul Keres
\end{frame}

\begin{frame}{Video Animation of optimal ply vs our policy}
  maybe for a mate in 5 or more

\end{frame}

\begin{frame}{Future Work}
  \begin{itemize}
    \item Zobrist hashing and invariance properties in chess endgames to reduce state space cardinality
    \item Implementing draw by repetition and 50-move rule
    \item Other chess endgames (e.g., KBBK, KRPK)
    \item Experiment with training against random player
  \end{itemize}
\end{frame}

{\setbeamercolor{palette primary}{fg=white, bg=bluscuro}
\begin{frame}[standout]
\thispagestyle{empty}
  {\LARGE Thank You!}
\end{frame}
}

\begin{frame}
  Why $\gamma=0.99$?
  
  The longest rkr endgame requires 16 white moves, hence $\gamma^{16} \approx 0.85$ is a good reward for reaching the goal and producing fast mates.
\end{frame}

